\chapter{Practical and Theoretical Rationality}
\label{sec:p11-rationality-2-2}

A belief about what will happen and a decision about what to do can concern
the same possible event while answering different evaluative questions.
Evidence that an action will occur need not make that action advisable;
conversely, the desirability of an outcome need not establish that it will
occur. This chapter reconstructs Ralph Wedgwood's account of the distinction
and of four attempts to explain what the two kinds of rationality have in
common. The selected source is his complete chapter~2.2 of
\emph{The Handbook of Rationality}, including its nine notes and references
\parencite[pp.~137--145]{RalphWedgwood2021}.

Wedgwood's investigation is philosophical. Historical interpretations,
proposed normative principles, and hypothetical cases are distinguished
below from mathematical consequences of explicit assumptions. The chapter
reports no participant study or estimate of how often people conform to a
standard. Its descriptions of contemporary philosophical agreement are
the author's assessments in 2021, not measured consensus. As
Section~\ref{sec:p11-method-boundary} emphasizes, a philosophical proposal
does not acquire theorem status merely by being expressed precisely.

\section{The distinction concerns attitudes, not simply subject matter}
\label{sec:p11-rat22-attitudes}

\subsection{The historical routes and their difficulty}

Wedgwood begins with Aristotle's account of activities involving reason.
In the teleological interpretation he presents, an activity is understood
through the function it ideally fulfills in human life. Contemplation or
study, \emph{the\={o}ria}, manifests theoretical wisdom, \emph{sophia}:
scientific knowledge, \emph{epist\={e}m\={e}}, rests on comprehension of first
principles, \emph{nous}, and contemplation exercises that understanding.
The relevant \emph{Nicomachean Ethics} locators include 1141a18 and
1177a26. Practical wisdom or prudence, \emph{phron\={e}sis}, is another
intellectual virtue, distinct also from technical skill, \emph{techn\={e}}
(1140a1--b29).

In the favorable case, practical wisdom manifests itself in action,
\emph{praxis}, that executes a decision, \emph{prohairesis}. A decision is
not merely a prediction: it is a desire to act here and now, reached from a
wish for an end regarded as good through deliberation about achieving it
(1094a5, 1111b26, 1112b15). Correct practical exercise combines understanding
of the good with an accurate perception of the situation, producing a
correct decision and action (1140b21). Wedgwood restricts this discussion to
Aristotle's narrower, reason-involving use of \emph{praxis}. His note~1
contrasts the broader usage at 1111a26, which includes voluntary behavior
by children and nonhuman animals and therefore does not itself imply the
exercise of reason.

Another Aristotelian characterization, at 1139a3--15, distinguishes the
study of what cannot be otherwise from deliberation about what can be
otherwise. Taken as a division between necessary and contingent truths,
it is inadequate for the present distinction: inquiry into ancient
history or paleontology can concern contingent facts without having a
practical purpose. This is Wedgwood's objection to that characterization,
not a claim that every aspect of Aristotle's account reduces to it.
His note~2 identifies Irwin's 1999 translation as the translation used
\parencite[pp.~137--138 and p.~144, nn.~1--2]{RalphWedgwood2021}.

Kant's formulations encounter a related problem when the distinction is
made between kinds of cognition or propositions. In the Logic edited by
J\"asche, practical cognition is an imperative expressing a possible free
action through which an end is realized; theoretical cognition concerns
what is, rather than what ought to be. The source locates this formulation
in Kant's collected writings, volume~9, p.~86. In the \emph{Groundwork},
volume~4, pp.~413--414, an imperative expresses an objective principle as
binding a will that need not already conform to it.

A hypothetical imperative presupposes an end and identifies necessary
means. Its weak rational constraint is that a rational agent cannot will
the end without also willing the necessary means. Yet the connection
between the means and the end can itself be described as a theoretical
causal proposition. Kant's later discussion calls such imperatives
technical and distinguishes their practical form from their theoretical
content (volume~20, pp.~196, 200). Wedgwood takes this difficulty to favor
a distinction between attitudes toward propositions rather than between
two disjoint classes of propositional content. His note~3 identifies the
1992, 2000 and 2012 translations used; those references do not amount to a
complete reconstruction of Kant's theory
\parencite[p.~138 and p.~144, n.~3]{RalphWedgwood2021}.

\begin{hbdefinition}{def:p11-rat22-attitude-types}{Wedgwood's attitude-based distinction}
\emph{Theoretical rationality} evaluates beliefs and belief-like or
belief-involving states and events. These include outright belief,
credence, inference, and belief revision.
\emph{Practical rationality} evaluates practical attitudes and their
exercise: intentions, decisions or choices, preferences, and intentional
actions. An intention is a state; deciding can be an event of forming or
revising an intention; carrying out the intention is an action.

The same propositional content can enter both assessments. The distinction
concerns the attitude and its evaluative role, not whether its subject is
necessary, contingent, morally important, or useful to an immediate project.
This is the approach Wedgwood endorses as the usual contemporary starting
point, subject to the qualifications that follow.
\end{hbdefinition}

The theoretical side need not consist only of outright belief. Wedgwood
names Joyce's proposal of precise and imprecise credences, Edgington's
conditional and unconditional credences, and Friedman's account of
suspension of judgment as a distinctive attitude. He does not establish
in this chapter that all these attitudes exist or reduce to one another:
the claim is that, if they exist and are rational, their rationality is
theoretical. Belief-involving events likewise include acquiring an attitude
where there was none, replacing it with another attitude to the same
proposition, and reaffirming an existing attitude.

On the practical side, theories disagree about the fundamental object of
assessment. Wedgwood reports Harman's emphasis on processes of changing
plans and Broome's emphasis on sets of mental states rather than direct
requirements on actions. Formal decision theories associated with Jeffrey
and Joyce often place preferences and their adjustment to beliefs at the
center. Connecting preferences to action needs a further premise: an agent
does not choose or intentionally perform a course of action when an
available alternative is strictly preferred. That is a proposed
preference--choice link, not a consequence of representing an ordering
by numbers alone
\parencite[pp.~138--139]{RalphWedgwood2021}.

\begin{hbexample}{def:p11-rat22-intention-belief}{A prediction can be justified without making its execution advisable}
Wedgwood identifies a complication for accounts, such as Setiya's, that
identify intending to do something with a special belief that one will do
it. Its distinctive causal tendency is to help bring about its own content,
rather than merely to register an independently occurring event.

Suppose an agent's evidence makes it clear that the agent will perform an
action, but performing it is plainly foolish or inadvisable. A belief that
the action will occur can then be theoretically rational while the
intention, understood on this identification as that special belief, is
not practically rational. The example holds apart evidential support for
occurrence and the practical assessment of execution. A theory identifying
these attitudes cannot distinguish the two evaluations merely by assigning
one to belief and the other to a different type of state; it needs a
further distinction. Wedgwood therefore presents the attitude-based
division as broad agreement with a genuine exception, not unanimity.
\end{hbexample}

The ancient emphasis on knowledge and decision offers continuity with
the attitude-based account. In particular, many contemporary philosophers
regard knowledge as a best case of full belief, while Aristotle connects
practical wisdom with decision and action. Kant's identification of the
will with practical reason (volume~4, p.~412) offers a further continuity:
a determination of the will appears to involve an intention to act.
These continuities do not remove the contemporary disagreements about
which states and processes are fundamental
\parencite[pp.~138--139]{RalphWedgwood2021}.

\section{What a unifying account would have to explain}
\label{sec:p11-rat22-unification}

Wedgwood identifies several parallels. Beliefs and choices can both be
called rational, reasonable, or justified; both can be rationally required
or permissible; and both admit distinctions between some, sufficient,
and overriding reasons. He suggests that these uses express a common
concept rather than mere verbal ambiguity. Rationality also admits
degrees on both sides. These parallels motivate the search for unification;
similar vocabulary by itself does not prove a shared underlying theory
\parencite[pp.~139--140]{RalphWedgwood2021}.

\begin{hbdefinition}{def:p11-rat22-requirement-basis}{A requirement and a rational way of meeting it}
Wedgwood proposes the following way to understand a rational requirement:
a condition is rationally required of an agent at a time just when that
condition is met in every relevantly available world in which the agent
is thinking as rationally as the agent can at that time.
The relevant availability and the standard of best possible rational
thinking must be specified in an application; this characterization does
not calculate them. His note~4 explicitly points to a different conception
in Broome's 2013 chapter~7.

Having an attitude that it is rational for one to have is different from
having that attitude rationally. The latter concerns how the attitude is
held or formed, not merely whether its content or selected option is
eligible under the standard. On the theoretical side, the distinction is
often expressed as propositional versus doxastic justification. Wedgwood
extends the same contrast to intentions and choices.
\end{hbdefinition}

His example is luck: an irrational thinker might happen to believe a
proposition for which adequate justification is available, or happen to
choose an option that is rationally eligible. That coincidence does not
establish that the thinker believed or chose rationally. This is also
different from happening to believe a truth or to produce a good outcome.
Eligibility, rational formation or maintenance, and external correctness
are three questions, not three names for one fact.

For Wedgwood, unification requires general principles about rationality
and an explanation of each domain by combining those principles with
principles specific to its attitudes. Principles about belief should
yield the specifically theoretical features; principles about intention,
choice, and action should yield the practical features. A common label
without these domain-specific bridges leaves the explanatory task undone
\parencite[p.~140 and p.~144, n.~4]{RalphWedgwood2021}.

\section{Four routes and their bridge premises}
\label{sec:p11-rat22-four-routes}

\subsection{Formal representation and pragmatic vulnerability}

Wedgwood's first route is the formal decision-theoretic approach. His
survey describes a common representation schema with three requirements:

\begin{enumerate}
\item Higher utility represents and reflects strict preference between
prospects, while equality of utilities represents and reflects indifference.
\item Higher probability represents and reflects greater confidence in a
proposition, while equality represents and reflects equal confidence.
\item The utility assigned to a prospect equals its expected utility under
the representing probability function.
\end{enumerate}

Here a \emph{prospect} is an alternative in the decision model; it is not
automatically a proposition believed or an already executed action. The
third requirement needs the outcome structure and utility interpretation
of the particular decision theory. Wedgwood's survey does not supply one
general outcome model or a single theorem covering all the versions it
mentions. His first clause does have an elementary mathematical
consequence, independent of that missing model
\parencite[p.~140]{RalphWedgwood2021}.

\begin{hbtheorem}{thm:p11-rat22-numerical-order}{proved=annex:p11.rat22.numerical-order}{Ordering consequences of an exact scalar representation}
Let \(U\) be a real-valued function on the relevant prospects. Suppose an
agent strictly prefers \(A\) to \(B\) exactly when \(U(A)>U(B)\), and is
indifferent between them exactly when \(U(A)=U(B)\). Here \(A,B,C\) range
over prospects, not formulas.

Every pair has exactly one of the following comparisons: strict preference
in one direction, strict preference in the other, or indifference.
Moreover, if the agent strictly prefers \(A\) to \(B\), and either strictly
prefers \(B\) to \(C\) or is indifferent between them, then the agent
strictly prefers \(A\) to \(C\). The conclusions include equal arguments
and require no finiteness assumption.
\end{hbtheorem}

The theorem does not establish that every rational agent has such a
representation. That normative assertion is an additional bridge premise.
The converse existence problem also requires conditions on the domain and
comparisons; the elementary implication cannot serve as its proof.

Wedgwood's account highlights Savage's representation theorem: under
specified preference conditions, preferences determine a unique
probability function and a utility function with unit and zero point left
free. He stresses that this construction does not require separate
assumptions about the agent's beliefs. But reducing theoretical to practical
rationality would need more than that mathematical representation. It
would require that rational preferences satisfy the theorem's conditions,
that the resulting probability function supplies the agent's beliefs,
and that probabilistic representability exhausts rational belief. The
identification of beliefs with what is implicit in preferences is the
controversial pragmatist premise. Wedgwood reports that later decision
theorists, including Joyce, seek to avoid commitment to it
\parencite[pp.~140--141]{RalphWedgwood2021}.

This is an attributed description of Savage's role in the proposed
reduction, not a statement of Savage's full domain, axioms, theorem, or
proof. The finite ordinal result in
Theorem~\ref{thm:p17-finite-representation} and the finite-prize lottery
result in Theorem~\ref{thm:p17-finite-vnm} are different presentations.
The latter assumes supplied lottery probabilities, independence, and
continuity. It therefore cannot replace Savage's construction of
subjective probabilities from preferences, nor can its proof establish
the disputed belief-identification premise.

A second formal route seeks a common pragmatic defect: violations of
theoretical or practical requirements can make an agent vulnerable to
a Dutch book or a money pump. Wedgwood presents the intended result as
acceptance of a set of bets or exchanges producing monetary loss regardless
of what happens, and points in note~5 to H\'ajek's discussion. His objection
is that vulnerability under a possible arrangement need not produce loss
in the agent's actual situation: the relevant bets might never be offered,
and the beliefs might remain useful in the practical situations encountered.
Thus a sure-loss theorem and its significance as a fundamental account of
rationality must be evaluated separately
\parencite[p.~141 and p.~144, n.~5]{RalphWedgwood2021}.

Theorem~\ref{thm:p17-finite-price-coherence} supplies an exact finite
contract result with its own prerequisites: common settlement, reversible
transactions at quoted prices, unrestricted finite real portfolios, and
no transaction charges. Definition~\ref{def:p17-finite-price-coherence}
does not claim that an agent accepts every allowed transaction. The theorem
therefore supports a conditional mathematical comparison, not the premise
that monetary vulnerability is sufficient to explain every irrational
belief. A money-pump argument also needs its particular assumptions about
sequential consent, payments, and persistence of preferences.

\subsection{Instrumental rationality and cognitive goals}

The second route makes instrumental rationality primary: rational thought
and action select means that appear optimal, given the evidence, for an
agent's goals. For theoretical rationality this would mean selecting
beliefs as means to a cognitive goal. A usual proposed goal is true belief,
without false belief, about the questions considered. On this proposal,
a belief set is rationally permissible exactly when it appears to be an
optimal means to that goal. This adds both a goal premise and a means--end
account of what makes belief rational
\parencite[p.~141]{RalphWedgwood2021}.

Wedgwood describes the approach as Hume-inspired, not as an uncontested
interpretation of Hume's \emph{Treatise}, II.iii.3. His note~6 records
Millgram's rival interpretation: Hume denied the strict application of
rationality predicates to actions and volitions, rather than advocating
the contemporary instrumental theory. Note~7 points to
Steglich-Petersen's 2018 treatment as a recent example of epistemic
instrumentalism. Neither historical ancestry nor the existence of that
example proves the proposed reduction.

\begin{hbexample}{def:p11-rat22-local-goals}{Flights, unwanted questions, and a goal-relative bridge}
In Kelly's objection as presented by Wedgwood, a traveler's interest in
accurate flight times from Denmark to California comes from an impending
choice of how to travel. The traveler has no comparable interest in flight
times between Cairo and Moscow because travel on that route is not
anticipated. These are local cognitive goals, not evidence that the agent
has adopted a general goal of believing every truth.

Now suppose an annoying interlocutor forces the agent to consider a
proposition \(p\), and compelling evidence for \(p\) is simultaneously
presented. The source judges that theoretical rationality requires belief
in \(p\), even if obtaining a true belief about it was not one of the
agent's goals. If this normative judgment is accepted, a theory that
requires rational belief only as a means to the agent's actual cognitive
goals needs an additional explanation of the requirement in this case.
The missing goal cannot be supplied by simply relabeling the evidence as
something the agent desires.
\end{hbexample}

The contrast is between individual, contingent interests and an
apparently intersubjective evidential requirement. It is not a population
study of travelers' goals. One may dispute the stipulated rational
requirement or offer a different account of cognitive goals; either reply
changes a substantive premise. Wedgwood calls continued investigation of
the reduction worthwhile while reporting low contemporary optimism about
it. This is his dated evaluation of the philosophical debate, not a
quantitative estimate of its prospects
\parencite[p.~141 and p.~144, nn.~6--7]{RalphWedgwood2021}.

\subsection{Reasons-centered unification}

The third route explains both domains through appropriate responsiveness
to reasons. Practical responsiveness concerns choices and the formation,
revision, and maintenance of intentions or plans; theoretical responsiveness
concerns the corresponding treatment of beliefs. But there are distinct
versions. Wedgwood reports Raz's account of irrationality as malfunction of
rational powers, understood as capacities to recognize and respond to
reasons, rather than simply identifying rationality with successful
response. Kolodny instead seeks a reduction to correct response to reasons
the agent believes to exist
\parencite[p.~142]{RalphWedgwood2021}.

\begin{hbdefinition}{def:p11-rat22-reason-propositions}{One interpretation of objective and subjective reasons}
Fix the propositions that count as reason-propositions for an agent at a
time. On the interpretation Wedgwood describes, the objective reasons are
those reason-propositions that are true, whereas the subjective or apparent
reasons are those that the agent believes. A response-to-apparent-reasons
account adds the premise that responding appropriately to the latter is
what rationality consists in.

This is one interpretation of the distinction. It does not establish that
every believed proposition counts as a reason-proposition, that subjective
reasons are true, or that merely believing a reason settles how it combines
with other reasons.
\end{hbdefinition}

Note~8 names Parfit's sufficient-apparent-reason account. The stronger
actual-reasons identity, the different beliefs-about-reasons variants,
and their dependence problems are treated in
Sections~\ref{sec:p11-reasons-identity} and~\ref{sec:p11-reasons-believed}.
Wedgwood expressly directs the reader to Broome's preceding source chapter
and to his 2013 chapters~5--6 for objections. Broome's refusal of a
reduction must not be erased by using the common phrase
\emph{responding to reasons} for several different proposals.

Wedgwood's further question is whether \emph{reason} itself is unified
enough to explain the two domains. In Raz's characterization, standard
epistemic reasons form part of the case for a proposition's truth; reasons
for action concern respects in which actions are good. A demon's credible
threat to punish an agent unless a belief is held can supply a
nonepistemic consideration, without thereby supplying a case for that
belief's truth. Raz distinguishes such a consideration from standard
epistemic reasons and denies that truth or true belief is a value.
The practical and epistemic roles therefore cannot simply be declared
identical.

Raz's proposed connection appeals to their guiding role in thought and
to dispositions constitutive of intelligible thinking and agency:
acting in ways believed good and believing on truth-related grounds.
The agency premise is disputed. Wedgwood cites Setiya's objections to
the claim that intentional action always occurs under a representation
of the good. A reasons-centered explanation accordingly owes either a
defense of this shared role or another bridge between the domains
\parencite[p.~142 and p.~144, n.~8]{RalphWedgwood2021}.

Schroeder's Hume-inspired version instead grounds reasons in desires or
similar attitudes. On the account Wedgwood surveys, a reason to perform an
act consists in a fact explaining how performing it promotes, or improves
the chances of, a desired state of affairs relative to the relevant
alternative. Extending this account to belief is a substantive proposal,
not a consequence of the practical formulation.

Wedgwood recounts two auxiliary claims meant to support that extension.
First, whatever an agent desires, refraining from a false belief would
at least somewhat improve the prospects of those desired outcomes,
yielding some reason against the false belief. Second, the weight of
a reason is not automatically the strength of the underlying desire or
the size of the improvement in its prospects. Schroeder supplies a
different account of weight elsewhere, which might allow the reasons
against false belief normally to be strongest. The first claim is not
proved here for every possible desire, and \emph{some reason} does not
entail \emph{most reason}. The second claim names an additional account
that must do substantive work; it is not itself that account. Wedgwood
therefore treats this version as a promising but unfinished project whose
ability to accommodate enough judgments about belief, choice, and action
remains to be established
\parencite[pp.~142--143]{RalphWedgwood2021}.

\subsection{Wedgwood's values-centered proposal}

The fourth route begins with rationality's ranking and guiding roles.
Some possible beliefs or choices are more rational than others, and
judgments about those rankings normally guide thinking toward the more
rational alternatives. Wedgwood argues that this resembles a value:
rational thinking is good in a distinctive respect. This is an argument
for treating rationality as a value, not a definition that settles the
issue without further defense
\parencite[p.~143]{RalphWedgwood2021}.

\begin{hbdefinition}{def:p11-rat22-correctness}{Rationality and correctness in the values-centered account}
Wedgwood's proposed account distinguishes an attitude's rationality from
its correctness or appropriateness. Rationality depends on its relations
to mental states and events present in the thinker at the relevant times.
Correctness depends on its relation to the external world.

For belief, the standard of correctness concerns whether the believed
proposition is true. For choice, it concerns whether the selected option
is feasible and genuinely good to do. These are different standards within
the proposed shared account. The mental-dependence premise and the
assessment of practical goodness are philosophical commitments, not
findings of the scalar-representation theorem.
\end{hbdefinition}

Correctness can also admit degrees. Choices can be more or less badly
wrong. Wedgwood proposes interpreting the inaccuracy measures discussed
by formal epistemologists such as Joyce as measures of how incorrect
credences are. For a true proposition \(p\), greater confidence in \(p\)
is more correct on this interpretation; for a false proposition, greater
confidence is less correct. An irrational belief can therefore be correct
by luck. This is a proposed interpretation of graded accuracy, not a
specific scoring rule, a guarantee of calibration, or a claim that
confidence alone reveals which proposition is true.

Expectations of correctness could guide the selection of beliefs and
choices. Thus correctness, like rationality, can rank alternative attitudes
and guide reasoning. Wedgwood takes this to support regarding correctness
as a value too. To connect the two, he supposes that, for each agent at a
time, there is a \emph{rational probability function} \(P\), and that
credences are rational in proportion to how closely they approximate it.
Relative to \(P\), greater irrationality would then mean a greater expected
shortfall from the highest available expected degree of correctness.
This is his proposal that irrationality provides adverse probabilistic
information about correctness
\parencite[p.~143]{RalphWedgwood2021}.

\begin{hbwarning}{def:p11-rat22-value-bridge}{The probability and value premises are not yet an implementation}
The selected chapter supposes the rational probability function; it does
not construct it, prove its existence or uniqueness, specify a metric for
approximation, choose an accuracy measure, or supply the available-attitude
domain needed to calculate the optimum. It also does not prove that
closeness to \(P\) has the proposed expected-correctness relationship.
Those are substantive dependencies of a developed theory, not consequences
of writing \(P\).

In particular, an agent's actual confidence, a preference-derived
probability, and a rational probability function must not be identified
without an argument. Nor does expected correctness guarantee actual
correctness on a particular occasion. Accepting a common expected-correctness
idea still leaves the distinct correctness standards for belief and choice
to be justified.
\end{hbwarning}

Raz's rejection of true belief as a value supplies an objection to this
route. Wedgwood responds that the relevant value is distinctive. An
admirable object exemplifies a value that guides one toward an attitude of
admiration. Correctness instead belongs to attitudes, including admiration
itself. Judging a belief or choice correct directly guides one toward
having that belief or making that choice; it does not first direct one to
adopt a further evaluative attitude toward the attitude. Wedgwood argues
that rationality judgments guide in the same direct way. This is meant
to make their value status intelligible without equating correctness with
the familiar value of an admirable object.

The unifying move is therefore neither to impose truth as the practical
standard nor to impose desired outcomes as the epistemic standard. It
applies one proposed connection between rationality and expected
correctness to two different types of attitude, with truth supplying the
belief standard and feasibility and goodness supplying the choice
standard. Wedgwood ends conditionally: a successful defense of this
account would explain rationality on both sides through its favorable
probabilistic connection to correctness. His note~9 directs the reader
to his 2017 book for the fuller values-centered theory; that theory is
not supplied by this survey chapter
\parencite[pp.~143--144 and p.~144, n.~9]{RalphWedgwood2021}.

\section{Applying the distinctions without crossing their boundaries}
\label{sec:p11-rat22-use}

\begin{hbmethod}{def:p11-rat22-bridge-check}{Checking a proposed belief--action bridge}
Identify what is evaluated: propositional content, a belief or credence,
an intention, a preference, a transition between attitudes, or an executed
action. State the information, alternatives, and relevant time. Do not
classify the case solely by whether its subject is practical.

Separate rational eligibility, the way the attitude was formed or held,
and external correctness. A lucky true belief or a lucky successful choice
need not establish rational formation; an eligible option need not have
been selected on an appropriate basis.

Specify the proposed unification and its added premises. A representation
argument needs its domain, comparison conditions, and any identification
of numerical probabilities with beliefs. An instrumental argument needs
the cognitive goals and the means--end criterion. A reasons argument
needs its objective, apparent, or believed-reason interpretation, its
account of weight, and a shared role that connects epistemic and practical
reasons. A values argument needs its correctness standards, probability
function, and justification of the expected-correctness bridge.

State separately what follows mathematically from the assumptions and
what is being proposed normatively or conceptually. Check the source's
contrasts: an evidenced prediction of foolish action, justified belief
without a desired cognitive goal, a threatened belief unsupported by
truth-related evidence, or a sure-loss portfolio never offered.

If the claim concerns actual people, supply the task, measure, population,
design, uncertainty, and competing explanations. Neither a philosophical
example nor a formal representation provides that empirical record.
\end{hbmethod}

This is an editorial use of Wedgwood's distinctions, not his proposed
decision algorithm. A successful application may locate a missing bridge
without selecting one philosophical theory. For example, information
supporting a traveler's prediction of a missed flight bears on theoretical
rationality; deciding whether to take that flight also requires alternatives,
feasibility, and a practical standard. Expected-utility calculations become
available only after the relevant decision model has been supplied. They
do not make the evidential prediction and practical evaluation the same
assessment.

The P17 finite theories establish precise representation and coherence
results under their stated mathematical premises. The present chapter
asks what, if anything, makes such premises fundamental requirements of
rational thought across attitude types. The Broome treatment in
Chapter~\ref{sec:p11-reasons-and-rationality} examines a different proposed
reduction and preserves the difference between correct response, coherence,
and achieved success. The P19 distinction between a benchmark and an
observed response, Definition~\ref{def:p19-benchmark}, explains why an
observed conformity rate could not by itself justify a benchmark;
Definition~\ref{def:p19-evidence-record} gives the record required for a
genuinely empirical claim. None of these interfaces turns one source's
normative or descriptive meaning of rationality into a universal house
definition.

The selected chapter supplies historical interpretation, an attitude
distinction, an explanatory target for unification, four routes with
objections, and a conditional values-centered proposal. Its references
identify further theories and objections rather than importing their
complete proofs or evidence into this treatment. The outstanding details
of those separate works remain separate. What is established locally is
the elementary ordering consequence in
Theorem~\ref{thm:p11-rat22-numerical-order}; what remains philosophical is
the choice and defense of the proposed bridge between belief and action
standards
\parencite[pp.~144--145]{RalphWedgwood2021}.

\index{rationality!practical and theoretical}
\index{rationality!instrumental}
\index{rationality!reasons-centered accounts}
\index{rationality!values-centered account}
\index{belief!rationality and correctness}
\index{choice!rationality and correctness}
